\documentclass[conference]{IEEEtran}
\usepackage{cite}
\usepackage{amsmath,amssymb,amsfonts}
\usepackage{algorithmic}
\usepackage{graphicx}
\usepackage{textcomp}
\usepackage{xcolor}

\begin{document}

\title{Algosee: A Machine Learning-Driven Framework for Explainable Algorithm Visualization}

\author{\IEEEauthorblockN{Anonymous Authors}
\IEEEauthorblockA{\textit{Department of Computer Science} \\
\textit{University/Institution}\\
City, Country \\
email@domain.com}}

\maketitle

\begin{abstract}
Traditional algorithm visualization tools often fail to demonstrate the impact of real-world data distributions on algorithmic efficiency, relying strictly on theoretical models that treat all inputs as uniform. This paper presents \textbf{Algosee}, a novel full-stack framework that integrates empirical benchmarking with a Random Forest-based prediction engine to provide data-specific algorithm recommendations. By extracting high-dimensional statistical features from non-uniform input arrays, Algosee bridges the gap between theoretical Big-O time complexity and practical runtime behavior. Furthermore, by employing Explainable AI (XAI) through Large Language Models (LLMs), the system translates complex execution metrics and statistical thresholds into pedagogical, human-readable insights. Evaluated on a diverse dataset of 60,000 algorithmic benchmarks, our predictive model achieves a \textbf{98.80\% accuracy} in determining the optimal sorting algorithm. We ultimately demonstrate how the fusion of machine learning prediction and generative AI explanation establishes a new pedagogical paradigm in continuous computer science education.
\end{abstract}

\begin{IEEEkeywords}
Algorithm Visualization, Machine Learning, Explainable AI, Educational Data Mining, Time Complexity
\end{IEEEkeywords}

\section{INTRODUCTION}
The analysis and optimization of sorting algorithms represent a fundamental pillar of computer science. While traditional algorithm analysis focuses primarily on theoretical asymptotic time and space complexity---expressed via Big-O, Big-Theta, and Big-Omega notation---real-world computational performance is significantly influenced by hardware-level factors and input data patterns. Theoretical analysis inherently assumes a uniform cost model; however, practical behavior is governed by branch prediction penalties, CPU cache locality, variable memory latency, and explicit input configurations such as randomness, partial ordering, duplication, and array dimensionality \cite{cormen_intro}.

Educational platforms and visualization tools have historically struggled to convey these practical nuances. Students are routinely taught that Quick Sort averages $O(N \log N)$ and Insertion Sort executes in $O(N^2)$, yet standard curricula often neglect to emphasize that for small arrays or highly pre-sorted arrays, Insertion Sort vastly outperforms Quick Sort due to constant-factor overheads and strict cache coherence.

Algosee addresses this structural deficit by aiming to bridge the gap between theoretical abstraction and empirical reality. By leveraging a machine learning-driven architecture, it combines explicit real-time algorithm visualization with empirical runtime benchmarking. Algosee extracts high-dimensional array features---such as Shannon entropy, inversion counts, and local sorted segments---and passes them through a Random Forest Classifier to dynamically predict the optimal execution path. Consequently, it enables students to understand not just \textit{how} an algorithm physically mutates variables, but \textit{why} a specific algorithmic selection represents the most mathematically optimal strategy for a localized data context. 

\section{RELATED WORK}
The evolution of Algorithm Visualization (AV) represents a gradual shift from static graph diagrams to highly interactive, web-based environments. Despite significant improvements in user interface design, a considerable gap persists in algorithmic personalization via data science.

\subsection{Foundational Algorithm Visualization}
Early empirical research into educational technologies, notably by Hundhausen et al. \cite{hundhausen2002meta}, established that the pedagogical efficacy of AV tools relies far more on the level of cognitive engagement required from the user than on the graphical fidelity of the animation itself. Platforms such as VisuAlgo \cite{halim2012visualgo} and Algorithm Visualizer have since become pervasive industry standards. These systems offer robust, interactive visualizations that execute state definitions across a multitude of languages. However, they remain fundamentally ``input-agnostic." They uniformly treat a deeply disordered high-entropy dataset identically to an array experiencing minor systemic noise, entirely failing to illustrate the adaptive triggers utilized by modern hyper-heuristics (e.g., Timsort's dynamic partitioning).

\subsection{Machine Learning in Educational Tools}
Recent methodological shifts toward Educational Data Mining (EDM) advocate for adaptive learning software. Historically, Breiman \cite{breiman2001random} demonstrated the profound robustness of Random Forest ensembles when subjected to high-dimensional, non-linear classification landscapes. While traditional machine learning in computer science education has predominantly focused on student dropout prediction or grading regression, Algosee uniquely subverts this paradigm. Algosee applies the ML ensemble directly to the computer science theorem itself, acting as a dynamic meta-classifier that predicts algorithmic efficiency on the fly based on array composition.

\subsection{Generative LLMs and Explainable AI (XAI)}
The accelerating emergence of Large Language Models (LLMs) has fundamentally redefined automated intelligent tutoring architectures. Recent investigations by Sarsa et al. \cite{sarsa2022automatic} emphasized the effectiveness of utilizing generative transformer models to compose conceptually accurate code explanations. Algosee extends this frontier by integrating modern, accelerated LLM inference APIs (such as Llama-3-70B via Groq) to transition from static software tooltips to dynamic Explicit Explainable AI (XAI). This integration successfully maps the ``black-box" numerical outputs of the classification engine to natural language rationales, delivering pedagogical context historically absent in AV software.

\section{PROPOSED METHODOLOGY}
The underlying Algosee architecture departs from conventional frontend visualizers by implementing a unified, full-stack ML pipeline. The architecture is systematically partitioned into three primary modules: the Feature Engineering Engine, the Predictive Optimization Module, and the Generative XAI Pipeline.

\subsection{Feature Extraction and Engineering}
Prior to execution or rendering, the system intercepts the user's input array $A$ of size $N$ to compile a topological and statistical signature. This deterministic engineering phase calculates features crucial for defining an array's systemic ``difficulty." 

The extracted numerical features encompass:
\begin{enumerate}
\item \textbf{Sortedness Score:} The ratio of properly ascending adjacent pairs over the total pairs $N-1$, providing a linear heuristic of ordered states.
\item \textbf{Inversion Count:} An absolute measure of global mathematical disorder, conceptually representing the minimum number of adjacent swaps required to perfectly sort the array, computed via a modified $O(N \log N)$ merge-sort algorithm.
\item \textbf{Shannon Entropy calculation:} Indicates the information density and duplicate probability of the value distribution, computed as $H(X) = -\sum_{i=1}^{n} P(x_i) \log_2 P(x_i)$. Lower entropy values signify numerous duplicates, a metric directly impacting algorithms like 3-Way Quick Sort.
\item \textbf{Variance \& Standard Deviation:} Evaluates the spatial dispersion and spread of the localized integers.
\item \textbf{Presorted Segments:} Identifies contiguous, naturally occurring sorted sub-sequences (Runs). High counts directly favor Insertion Sort variants.
\end{enumerate}

\subsection{Predictive Modeling and Dataset Generation}
To train a robust heuristic engine, a massive synthetic dataset of $60,000$ discrete integer arrays was procedurally generated. These arrays ranged in uniform size constraints ($10 \leq N \leq 5000$) and were subjected to various distribution alterations inclusive of absolute random, nearly-sorted (Gaussian noise application), reversed, and heavily duplicated contexts. 

Through exhaustive grid-search hyperparameter tuning, a \textbf{Random Forest Classifier} ($n_{estimators}=200, \text{criterion}='gini'$) was empirically selected as the optimal model. Random forests inherently provide substantial resilience against multi-collinearity among array statistical features and output critical feature-importance scores natively. Deep Neural Networks (DNNs) were explicitly bypassed for this implementation due to their lack of deterministic explainability and high latency overhead when parsing nominal array lengths. 

During runtime, the Random Forest Classifier maps the instantaneous feature array into a probability distribution predicting the lowest absolute physical execution latency, successfully classifying the ``Winner" algorithm out of six classical configurations (Bubble, Selection, Insertion, Merge, Quick, Heap).

\subsection{Explainable AI (XAI) Integration}
To translate the purely mathematical predictions into actionable tutoring interventions, Algosee utilizes few-shot and zero-shot prompt engineering against the Llama-3 meta-heuristic APIs. When the Random Forest outputs a prediction, the backend constructs a highly structured JSON query injecting the raw metric data (\texttt{swaps}, \texttt{comparisons}, \texttt{inversions}, \texttt{predicted winner}). The transformer returns heavily categorized reasoning:
\begin{itemize}
\item \textbf{Conceptual Overviews:} Provides the fundamental algorithmic mechanics.
\item \textbf{Performance Verdicts:} Synthesizes the exact reason the Random Forest predicted a specific winner (e.g., ``Insertion Sort was recommended because the array recorded an exceptionally low inversion count of 4").
\item \textbf{Optimization Advice:} Curates professional, data-specific engineering techniques and memory tuning recommendations.
\end{itemize}

\section{EXPERIMENTAL RESULTS \& ANALYSIS}
To ensure the statistical validity of the classification framework, the Random Forest model was rigorously evaluated utilizing a sequestered, balanced test set containing 2,000 randomized array environments.

\subsection{Overall Performance Metrics}
The predictive engine within Algosee successfully recorded a \textbf{Total Validation Accuracy of 98.80\%}. The confusion matrix distribution effectively highlights an absence of overt classification biases. A summary of the per-algorithm predictive metrics across standard benchmarks is documented below:

\begin{table}[h]
\centering
\caption{Overall Performance Metrics}
\label{tab:performance_metrics}
\begin{tabular}{|l|c|c|c|c|}
\hline
\textbf{Algorithm} & \textbf{Precision} & \textbf{Recall} & \textbf{F1-Score} & \textbf{Latency} \\
\hline
\textbf{Insertion Sort} & 99.1\% & 98.7\% & 98.9\% & $<$ 5 ms \\
\textbf{Quick Sort} & 99.4\% & 98.8\% & 99.1\% & $<$ 5 ms \\
\textbf{Merge Sort} & 97.9\% & 99.6\% & 98.7\% & $<$ 5 ms \\
\textbf{Heap Sort} & 86.3\% & 79.1\% & 82.6\% & $<$ 5 ms \\
\hline
\end{tabular}
\end{table}

\subsection{Analytical Discussion and Feature Importance}
Observations drawn from the experimental outcomes denote exceptional precision capabilities in segregating traditional threshold barriers. Through Gini importance analysis of the Random Forest, \textbf{Inversion Count} and \textbf{Array Size ($N$)} were revealed to be the dominant decision nodes. The model demonstrated profound competency at identifying the $N < 50$ boundary layer---accurately transitioning recommendations from Quick Sort to native Insertion Sort in direct reflection of actual CPU cache constraints.

Conversely, the moderately reduced F1-Score relating to Heap Sort (82.6\%) traces its origin to inherent theoretical overlaps. Both Merge Sort and Heap Sort strictly adhere to an $O(N \log N)$ overarching time guarantee under nearly all input distributions. As a direct result, absolute runtime discrepancies between the two are frequently governed by minute physical CPU architectural factors or contextual operating system thread management, rather than the explicit array topology. Despite this microscopic noise, the overarching heuristic remains highly dominant against naive, unified sorting methodologies.

\subsection{Case Study: High Entropy versus Near-Sorting}
In systematic testing, an array initialization containing 10,000 highly duplicated elements (Entropy $H \approx 0.15$) triggered the LLM subsystem to recommend explicitly against utilizing naive Quick Sort, directly cautioning the user regarding $O(N^2)$ worst-case deterioration without 3-way partitioning. Alternatively, supplying a list of sequentially ascending integers with a singular misplaced digit ($N=100$) invoked an instantaneous 99\% probability prediction for Insertion Sort, demonstrating a clear mathematical comprehension of $O(N)$ best-case triggers.

\section{CONCLUSION \& FUTURE WORK}
This research delineates the architecture and deployment of Algosee, a machine learning-augmented visualization platform uniquely positioned at the intersection of conventional pedagogical animation and explicit data science. By integrating a heavily optimized Random Forest classifier with dynamic LLM-driven Explainable AI generation, the framework effectively breaks the constraints of uniform, input-agnostic educational software. It exposes the hidden reality of data engineering: that theoretical asymptotic complexity curves rarely reflect true physical operational limits without context.

Future expansion of the Algosee framework proposes extending empirical modeling capabilities beyond strict linear sorting schemas to encompass computationally intensive graph traversal permutations, dynamic programming optimizations, and localized memory-allocation constraints. Further integration of reinforcement learning (RL) models could potentially allow the system to self-optimize and auto-generate specialized hybrid algorithms uniquely constructed for highly anomalous datasets.

\begin{thebibliography}{1}

\bibitem{cormen_intro}
T.~Cormen, C.~Leiserson, R.~Rivest, and C.~Stein, \emph{Introduction to Algorithms}, 3rd~ed.\hskip 1em plus 0.5em minus 0.4em\relax MIT Press, 2009.

\bibitem{hundhausen2002meta}
C.~D. Hundhausen, S.~A. Douglas, and J.~T. Stasko, ``A Meta-Study of Algorithm Visualization Effectiveness,'' \emph{Journal of Visual Languages \& Computing}, vol.~13, no.~3, pp. 259--290, 2002.

\bibitem{halim2012visualgo}
F.~Halim, Z.~C. Koh, V.~B. Loh, and S.~Halim, ``VisuAlgo: Visualising Data Structures and Algorithms through Browsers,'' \emph{IEEE Transactions on Learning Technologies}, 2012.

\bibitem{breiman2001random}
L.~Breiman, ``Random Forests,'' \emph{Machine Learning}, vol.~45, no.~1, pp. 5--32, 2001.

\bibitem{sarsa2022automatic}
S.~Sarsa, P.~Denny, A.~Hellas, and J.~Leinonen, ``Automatic Generation of Programming Exercises and Code Explanations Using Large Language Models,'' in \emph{Proceedings of the 2022 ACM Conference on International Computing Education Research (ICER)}, 2022.

\bibitem{chawla2002smote}
N.~V. Chawla, K.~W. Bowyer, L.~O. Hall, and W.~P. Kegelmeyer, ``SMOTE: Synthetic Minority Over-sampling Technique,'' \emph{Journal of Artificial Intelligence Research}, vol.~16, pp. 321--357, 2002.

\end{thebibliography}

\end{document}
